\chapter{Нечётное продолжение хопфовой пары и равногранный стоп-тест}

Литературный аудит заменил образ ``линия проходит через ядро'' на более
точную конструкцию: нечётный оператор $T(x)$ обратим вне дефекта и теряет
ранг на дефекте. Настоящий гейт проверяет все равногранные реализации,
доступные текущему родителю, и отдельно перепроверяет, не является ли
разность рангов $20-15$ полезным относительным индексом вместо
препятствия.

\section{Прямой прямоугольный оператор}

Естественная нечётная компонента связывающей алгебры имеет тип
\begin{equation}
 T:\mathbb C^{15}\longrightarrow\mathbb C^{20}.
 \label{eq:v5-odd-core-rectangular-T}
\end{equation}
Соответствующий самосопряжённый оператор равен
\begin{equation}
 Q_T=
 \begin{pmatrix}
 0&T\\T^*&0
 \end{pmatrix}
 \quad\text{на}\quad
 \mathbb C^{20}\oplus\mathbb C^{15}.
 \label{eq:v5-odd-core-selfadjoint-Q}
\end{equation}

Если $T$ имеет максимальный ранг пятнадцать, то
\begin{equation}
 \dim\ker T=0,
 \qquad
 \dim\ker T^*=5,
 \qquad
 \dim\ker Q_T=5.
 \label{eq:v5-odd-core-asymptotic-nullity}
\end{equation}
Следовательно, даже вдали от ядра полная матричная масса имеет пять
нулевых каналов. Это нарушает ключевое условие индекса Каллиаса:
асимптотическая матричная масса должна быть обратима, или по крайней мере
иметь выведенное равномерно обратимое ограничение. В стандартной формуле
масса на сфере бесконечности деформируется в $GL(N,\mathbb C)$; см.
\path{arXiv:2106.01591}, раздел 4.3.2.

При падении ранга на единицу получаем
\begin{equation}
 \dim\ker T=1,
 \qquad
 \dim\ker T^*=6,
 \qquad
 \dim\ker Q_T=7.
\end{equation}
Число нулей увеличивается на два, но прямоугольный индекс не меняется:
\begin{equation}
 \operatorname{ind}T
 =\dim\ker T-\dim\ker T^*
 =15-20=-5.
 \label{eq:v5-odd-core-constant-index}
\end{equation}
То есть локальная потеря одного сингулярного значения не создаёт индекс
один; она добавляет сопряжённую пару нулей поверх постоянного индекса
$-5$.

\begin{nogocriterion}[Пять асимптотических нулевых каналов]
Прямой нечётный угол $20\times15$ не задаёт оператор Каллиаса на полном
пакете: $Q_T$ не имеет щели на бесконечности, а локальное падение ранга не
меняет его постоянный прямоугольный индекс $-5$.
\end{nogocriterion}

\section{Структурный остаток \texorpdfstring{$1/7$}{1/7}}

Разность рангов не является произвольной. Для связывающей градуировки
\begin{equation}
 \Gamma_{\rm link}=p_{20}-p_{15}
\end{equation}
нормированный суперслед единицы равен
\begin{equation}
 \tau_{35}(\Gamma_{\rm link})
 =\frac{20-15}{35}=\frac17.
 \label{eq:v5-odd-core-one-seventh-superdimension}
\end{equation}
Это новый точный способ прочитать уже существующее число $1/7$: как
нормированную суперразмерность двух углов.

Но данный факт нельзя отождествлять с локализованной нейтринной модой или
физической массой. Суперразмерность является постоянным глобальным
инвариантом контейнера, тогда как нужный дефектный индекс должен возникать
из изменения асимптотического класса при сохранении щели.

\begin{remark}
Совпадение с прежним весом единственной голономной линии $1/7$ является
структурной зацепкой, но пока между двумя появлениями числа нет выведенного
операторного равенства.
\end{remark}

\section{Попытка квадратного угла
\texorpdfstring{$15\times15$}{15 x 15}}

Первый возможный ремонт --- выбрать ранг-пятнадцать подпространство в
левом углу $H_{20}$. Но замороженный семейный KO6-модуль состоит из
сопряжённой пары синглетов и трёх сопряжённых пар триплетов. Ранги
канонических редукций, построенных из этих целых блоков, имеют вид
\begin{equation}
 \{0,2,6,8,12,14,18,20\}.
 \label{eq:v5-odd-core-compatible-ranks}
\end{equation}
Ранг пятнадцать отсутствует.

Ранее выведенный опорный проектор имеет ранг два, а физический проектор ---
ранг восемнадцать. Получить пятнадцать можно лишь дополнительно удалить
три направления из физического угла. Ни алгебра, ни $J$, ни высота, ни
хопфов функтор не выбирают эти три направления.

Кроме того, такой квадрат $15\times15$ перестал бы быть прямым углом
$M_{20\times15}$ и потребовал бы нового условного следа. Поэтому он не
является скрытой частью текущего $M_{35}$.

\section{Попытка на паре
\texorpdfstring{$E\oplus E^*$}{E и E*}}

Пространства стрелок имеют равные комплексные размерности:
\begin{equation}
 E=M_{20\times15}(\mathbb C),
 \qquad
 E^*=M_{15\times20}(\mathbb C),
 \qquad
 \dim E=\dim E^*=300.
\end{equation}
На первый взгляд это даёт квадратный нечётный оператор ранга $300$.

Однако канонический переход
\begin{equation}
 X\longmapsto X^*
\end{equation}
антилинеен и меняет левое и правое действия местами. Комплексно-линейный
бимодульный изоморфизм $E\to E^*$ потребовал бы отождествить факторные
алгебры $M_{20}$ и $M_{15}$, что невозможно из-за различных рангов.

Если же объявить $E\oplus E^*$ новым пространством состояний, носитель
меняется с $300$ на $600$. В связывающей алгебре оба угла уже присутствуют
как операторные координаты, но это не означает, что оба являются
независимыми физическими состояниями. Такое чтение скрыто удваивает меру.

\begin{nogocriterion}[Антилинейность обратной стрелки]
Операция $*$ обеспечивает KO6-сопряжение $E\leftrightarrow E^*$, но не
даёт комплексно-линейное нечётное поле между двумя ранг-300 пакетами.
Превращение операторной пары в носитель ранга 600 является новой
архитектурой.
\end{nogocriterion}

\section{Самопроверка: квадрат
\texorpdfstring{$150\times150$}{150 x 150}}

Полный семейный модуль имеет две сопряжённые половины ранга десять. После
умножения на $H_{15}$ переходный носитель формально раскладывается как
\begin{equation}
 E=(H_{10}\otimes H_{15})
 \oplus J(H_{10}\otimes H_{15}),
 \qquad 300=150+150.
 \label{eq:v5-odd-core-150-pair}
\end{equation}
Эту возможность нельзя было исключать без проверки.

Но единственное каноническое отождествление двух половин снова задаётся
антилинейным $J$. Комплексно-линейное смешивание частицы с сопряжённой
частицей является майорановским типом спаривания. Пакет $H_{15}$ не
содержит стерильного калибровочного синглета $\nu_R$; он состоит из
$Q_L,L_L,u_R,d_R,e_R$. Поэтому универсальный обратимый оператор между
всеми 150 каналами нарушил бы квантовые числа Стандартной модели.

Хиггс может связать отдельные левые и правые блоки после нарушения
электрослабой симметрии, но прежние гейты уже показали, что эти компоненты
не образуют единственный полный оператор и оставляют независимые
юкавские направления. Следовательно, квадрат $150\times150$ также не
выводится.

\section{Почему положительные сжатия не сохраняют ориентацию}

Из прямоугольного $T$ автоматически получаются квадратные операторы
\begin{equation}
 T^*T\in M_{15},
 \qquad
 TT^*\in M_{20}.
\end{equation}
Они хранят сингулярные значения и локусы падения ранга. Но оба являются
чётными положительными операторами. Полярная фаза $T$, в которой могла бы
жить хопфова ориентация, при таком сжатии теряется. Поэтому $T^*T$ не
является нечётным отображением $L\leftrightarrow L^*$ и не несёт сам по
себе нужный класс.

\section{Можно ли вычесть постоянное коядро как фон}

В относительной $K$-теории допустимо сравнивать пакеты разных рангов и
рассматривать постоянную разность как фон. Тогда пять нулевых каналов
можно формально вынести из дефектного сектора.

Физически этого недостаточно. Пока не выведен проектор, отделяющий эти
пять каналов от оператора Дирака, они остаются безмассовыми состояниями на
бесконечности и уничтожают условие Фредгольма. Объявить их
``ненаблюдаемым фоном'' после расчёта означало бы вручную изменить
физический фактор.

Канонический проектор ранга пять в $H_{20}$ не найден: известные
совместимые редукции перечислены в
\eqref{eq:v5-odd-core-compatible-ranks}. Поэтому относительное вычитание
сейчас является математической перезаписью, а не родительским выводом.

\section{Итог исчерпания внутренней ветви}

Проверены четыре возможных равногранных чтения:
\begin{enumerate}
 \item прямой $20\times15$: пять постоянных асимптотических нулей;
 \item условный $15\times15$: нет канонического ранга-пятнадцать угла;
 \item стрелка/обратная стрелка $300\times300$: антилинейность и удвоение;
 \item частица/сопряжённая частица $150\times150$: нет универсального
 калибровочно-инвариантного линейного спаривания.
\end{enumerate}

Ни один вариант не даёт на текущем носителе нечётный оператор, который
одновременно:
\begin{equation}
 \begin{gathered}
 \text{обратим на сфере бесконечности},\\
 \text{теряет ровно один ранг на ядре},\\
 \text{сохраняет KO6, }H_{15}\text{ и след }M_{35}.
 \end{gathered}
\end{equation}

\begin{theorem}[Исчерпание нечётного хопфова продолжения версии V]
Хопфова линия и ориентационный функтор остаются корректными граничными
результатами. Однако текущий $H_{15}/M_{35}$ не содержит выведенного
равногранного нечётного оператора с асимптотической щелью и единичным
дефектным индексом. Прямой угол имеет постоянный индекс $-5$ и пять
нулевых каналов; все квадратные ремонты требуют ручной редукции,
комплексно-линейного удвоения или нового физического спаривания.
\end{theorem}

\section{Вердикт}

Внутренняя хопфова ветвь версии V исчерпана. Это не отменяет:
\begin{enumerate}
 \item хопфову линию $c_1=1$ на границе;
 \item функториальное назначение $L/L^*$ стрелкам $E/E^*$;
 \item точную суперразмерность $1/7$;
 \item существование внешнего математического BPS-завершения.
\end{enumerate}
Но ни один из этих результатов не создаёт физическую частицу на текущем
носителе.

Следующий документ должен быть не новой спасательной гипотезой, а
финальной заморозкой ретроспективного переоткрытия с условиями, при которых
оно может быть возобновлено уже в новой версии:
\path{version5_hopf_reopening_final_status_freeze_gate.tex}.

\begin{protocol}
\begin{enumerate}
 \item прямой нечётный угол: \textbf{$20\times15$};
 \item асимптотическая нульность $Q_T$: \textbf{пять};
 \item индекс $T$ при максимальном ранге: \textbf{$-5$};
 \item индекс после падения ранга на один: \textbf{$-5$};
 \item суперразмерность углов: \textbf{$1/7$};
 \item суперразмерность равна локальному индексу один: \textbf{нет};
 \item канонический угол $15\times15$: \textbf{нет};
 \item канонический линейный оператор $E\to E^*$: \textbf{нет};
 \item вариант $E\oplus E^*$ без удвоения: \textbf{нет};
 \item универсальный квадрат $150\times150$: \textbf{нет};
 \item относительное удаление пяти каналов выведено: \textbf{нет};
 \item граничная хопфова ориентация: \textbf{сохранена};
 \item физическое замыкание: \textbf{не достигнуто};
 \item внутренняя ветвь версии V: \textbf{исчерпана}.
\end{enumerate}
\end{protocol}

Машинный сертификат сохранён в
\path{s2t_v5_hopf_pair_odd_core_extension_gate_results.json}.